\documentclass[10pt]{article}
\usepackage[margin=0.8in]{geometry}
\usepackage{amsmath,amssymb,amsthm,mathtools,booktabs,microtype,hyperref}
\hypersetup{hidelinks,pdftitle={A canonical symmetry excludes the ramified quartic completion}}
\newtheorem{theorem}{Theorem}
\newtheorem{lemma}[theorem]{Lemma}
\newcommand{\R}{\mathbb R}
\newcommand{\Cr}{\mathcal C_{\mathrm{ram}}}
\newcommand{\Cc}{\mathcal C_{\mathrm{can}}}
\newcommand{\tr}{\operatorname{tr}}
\newcommand{\ip}[2]{\langle #1,#2\rangle}
\title{A canonical symmetry excludes the ramified quartic completion}
\author{}\date{}
\begin{document}
\maketitle\vspace{-2.5em}
\begin{abstract}
The fixed quartic selector is not universally adequate. A physical
Alice--Bob transport of an already certified source lies outside its cone.
Four explicit positive-definite integer block Hankel matrices separate this
all-positive rational source. The transport also identifies the fractional
Alice dependence missing from the present testing space. No larger family
is asserted to cover the full canonical cone.
\end{abstract}

\section{The counterexample and its validity}

Retain the accepted source order
$c=(d,\lambda,n,m,C,A,D,\alpha,\beta)\in\R_+^9$ and the complete
canonical parent $F_c$. The cone $\Cc$ consists of sources valid at every
state and every pair of projective measurement angles. The accepted
two-completion matrices are
\[
 G_J=\begin{pmatrix}P+K&J\\J^T&N+K\end{pmatrix},\qquad
 H_J=\begin{pmatrix}P+E&K-E-J\\K-E-J^T&N+E\end{pmatrix}.
\]
Here $P,N,E,K$ are exactly the source-affine quadratic matrices in the
canonical completion note~\cite{accepted}. The fixed cone $\Cr$ requires
$v=16s^2$, $s\ge0$, and one matrix polynomial
$J(s)=\sum_{k=0}^4J_ks^k$ making both matrices positive semidefinite.
Its order-49 Gram pencil is retained unchanged.

The established input used below is the exact admission
\[
 \bar c=\left(1,2^{-34},2^{-14},2^{-24},\frac1{576},161,
                   \frac12,76,\frac1{250}\right)\in\Cr\subseteq\Cc.
 \tag{1}
\]
Its existing proof and matrices are not recomputed.

\begin{theorem}\label{thm:main}
The all-positive rational source
\[
 \boxed{c^\dagger=\left(1,\frac1{576},161,2^{-24},2^{-34},2^{-14},
                   \frac12,\frac1{250},76\right)}           \tag{2}
\]
belongs to $\Cc\setminus\Cr$. Consequently
$\Cr\subsetneq\Cc$: not every valid canonical source admits the fixed
quartic selector. The accepted order-392 eight-copy outer lift is strictly
larger than $Q_{222}$, by the established polar assembly theorem~\cite{assembly}.
\end{theorem}

Define the involution
\[
 \pi c=(d,C,A,m,\lambda,n,D,\beta,\alpha),\qquad
 S(W,X,Y,Z)=(W,Y,X,Z).
\]
The full parent satisfies the exact identity
\[
 \boxed{u^2v^2F_{\pi c}(-1/v,-1/u;S\psi)=F_c(u,v;\psi).} \tag{3}
\]
For a direct verification, the eight event packets, in their accepted
order $(\lambda,n,m,C,A,D,\alpha,\beta)$, are carried respectively to
$(C,A,m,\lambda,n,D,\beta,\alpha)$. The $m,D$ roots change sign.
The $d$ coefficient is unchanged: its three positive packets are permuted
or fixed, and $D_uD_vZ^2$ is unchanged after clearing the angles.
This proves (3) coefficient by coefficient, without a determinant argument.

There is no omitted endpoint. With homogeneous angle pairs
$a=(a_0,a_1)$, $b=(b_0,b_1)$, identity (3) is the polynomial identity
\[
 \boxed{F^h_{\pi c}\bigl((b_1,-b_0),(a_1,-a_0);S\psi\bigr)
                       =F^h_c(a,b;\psi).}                 \tag{4}
\]
It holds on the complete product of angle circles, including literal zero
coordinates. It permutes zero source weights exactly, keeps all state
signs, and acts separately on source and testing coordinates. Real and
imaginary state parts give the same conclusion for complex states.
Thus $\pi\Cc=\Cc$. Since $c^\dagger=\pi\bar c$, (1) proves validity
of (2). It remains to exclude its quartic completion.

\section{An exact separating trace identity}

For $\sigma=G,H$, let $N_{\sigma,0},\ldots,N_{\sigma,4}$ be the ten
\emph{integer} symmetric $4\times4$ matrices in Appendix~\ref{app:data}.
Form the four integer block Hankel matrices
\[
 \mathsf U_\sigma=(N_{\sigma,i+j})_{i,j=0}^{2},\qquad
 \mathsf V_\sigma=(N_{\sigma,i+j+1})_{i,j=0}^{1}.           \tag{5}
\]
They have orders $12,8,12,8$ and are positive definite. Moreover,
\[
 (N_{G,k})_{\{1,2\},\{3,4\}}
 =(N_{H,k})_{\{1,2\},\{3,4\}}
 \qquad(0\le k\le4).                                    \tag{6}
\]
These properties are certified by the displayed integers and the exact
elimination recurrence in the appendix.

Write $Q(s)=\sum_{k=0}^4Q_ks^k$ for either completed matrix. Its accepted
half-line Gram representation is
\[
 Q(s)=Z_2(s)^TA_QZ_2(s)+sZ_1(s)^TB_QZ_1(s),\quad
 Z_2=(I_4,sI_4,s^2I_4)^T,\quad Z_1=(I_4,sI_4)^T,
\]
where $A_Q,B_Q\succeq0$. Expanding gives
\[
 \sum_{k=0}^4\ip{N_{\sigma,k}}{Q_k}
 =\ip{\mathsf U_\sigma}{A_Q}+\ip{\mathsf V_\sigma}{B_Q}.
                                                               \tag{7}
\]
Each $J_k$ occurs in opposite off-diagonal blocks of $G_k$ and $H_k$,
so (6) cancels every selector coefficient. Consequently the sum of (7)
is source-affine and nonnegative for every feasible certificate, including
zero weights and singular Grams.

Here is the resulting integer functional in the accepted source order:
\[
\begin{array}{c|r@{\qquad}c|r@{\qquad}c|r}
 d&-97633735551856&\lambda&24154193783510329&n&85793705264\\
 m&1787034292680750290&C&3303160925893294&A&16956553777510394\\
 D&13682527570574&\alpha&670221724556174&\beta&366431361010
\end{array}
                                                               \tag{8}
\]
Denote the functional with these coefficients by $\ell(c)$.
Direct substitution of the accepted $P,N,E,K$, with $v=16s^2$, gives
the exact identity
\[
 \boxed{12\ell(c)=
 \ip{\mathsf U_G}{A_G}+\ip{\mathsf V_G}{B_G}
 +\ip{\mathsf U_H}{A_H}+\ip{\mathsf V_H}{B_H}.}             \tag{9}
\]
For checking its nine coefficients, only the moments of degrees $0,2,4$
are needed. If $Q_j^{(v)}(c,0)$ is the coefficient of $v^j$ with $J=0$,
the left side of (9) is
\[
 \sum_{j=0}^2 16^j\bigl[
 \ip{N_{G,2j}}{G_j^{(v)}(c,0)}+
 \ip{N_{H,2j}}{H_j^{(v)}(c,0)}\bigr].                     \tag{10}
\]
Thus $\ell\ge0$ on $\Cr$. But direct rational evaluation yields
\[
 \boxed{\ell(c^\dagger)=
 -\frac{32607070960447356138172589}{9663676416000}<0.}      \tag{11}
\]
This proves Theorem~\ref{thm:main}. No closure or dual-attainment argument
is needed for this finite separating identity.

\section{The contact mechanism and the remaining problem}

The separation already holds at the exact contact source $\pi c_*$
underlying (1). Its contact in positive angle gauges is
\[
 u=\frac1{16},\qquad v=13,\qquad
 \psi=\left(140,-\frac1{50},\frac1{32},1\right),
\]
and
\[
 \ell(\pi c_*)=
 -\frac{4844338143654932884745484309348314410}
        {802811514300974151343419}<0.                    \tag{12}
\]
For clarity, both state-normal constraints are retained. Write
\[
 \eta=L(1/16)\psi=
 \left(\frac{71681}{8192},\frac{17}{6400},
       \frac{279}{32},-\frac{801}{800}\right).
\]
At $U=1/256$ and $s=\sqrt{13}/4$, any putative completion must satisfy
\[
 G(s)(\eta_1,\eta_2,U\eta_3,U\eta_4)^T=0,\qquad
 H(s)\eta^T=0.                                          \tag{13}
\]
These follow from the accepted split identity and positivity of both
summands at a zero. The separator excludes such a global quartic
completion and remains negative after opening both zero event weights to
the all-positive source (2).

The original proof's symmetry reduction preserves canonical validity,
but it does not preserve the present fixed testing space. Its five compact
Bob weights in affine coordinates are
\[
 \frac{(1,\sqrt{|v|}/4,|v|/16,|v|^{1/4}/2,|v|^{3/4}/8)}
      {\sqrt{1+v^2}}.
\]
Under (3), they become the fixed Alice weights
\[
 \boxed{
 \frac{(|u|,\sqrt{|u|}/4,1/16,|u|^{3/4}/2,|u|^{1/4}/8)}
      {\sqrt{1+u^2}}.}                                  \tag{14}
\]
Their branches and normalization depend only on the test. The homogeneous
pullback of the accepted forty functions represents $b_0^2F^h_{\pi c}$,
where the original family represented $a_1^2F^h_c$; both multipliers are
nonnegative and positive on a dense testing chart.

The new dependence is genuinely outside the old span. Restrict to
$u>0$, $v=1$, $\psi=(1,0,0,0)$. Every old compact feature is a constant
times a polynomial of degree at most three divided by $(1+u^2)^{3/2}$,
and is analytic at $u=0$. One transported feature is
\[
 -\frac{u^{1/4}}{16\sqrt2\sqrt{1+u^2}}.
\]
It has nonintegral order of vanishing and cannot lie in that span.
Over the common multiplier $a_1^2b_0^2$, multiply old roots by $b_0$ and
transported roots by $a_1$; the same feature becomes
$-u^{5/4}/(16\sqrt2(1+u^2))$, so the distinction persists.

In particular $\Cr$ and $\pi\Cr$ are incomparable: (1), (2), and
$\pi^2=1$ give a point in each set difference. Formula (14) is an exact
transport of an entire certificate cone, but it does not prove that its
addition covers $\Cc$. Nor does separation prove that a particular
quarter-power feature is the only possible repair.

\paragraph{Precise remaining implication.}
The implication $\Cc\subseteq\Cr$ is false, rather than unresolved.
The remaining task is to specify one finite, source-independent algebraic
testing space $\mathcal U$ and prove that every $c\in\Cc$ has its full
unmultiplied slack, or one fixed testing-multiplied slack, in the Gram cone
of $\mathcal U$. The R21 contact/conductor identities and qualitative
tensor criterion~\cite{companion} do not themselves identify that space.
No such universal conversion is established here. The accepted companion,
local certificates, and assembly results are unchanged.

\paragraph{Reproduction.}
Run \texttt{python3 verify.py}. The checker uses exact integer and rational
arithmetic, verifies (4) for all nine coefficient packets, verifies (5)--(11)
directly from the block definitions, and independently contracts the same
certificate against all 148 auxiliary and nine source matrices of the
unchanged public pencil. It also checks the value and all six first
derivatives at the contact in (13). No numerical optimizer is required.

\begin{thebibliography}{3}\small
\bibitem{accepted} \emph{Canonical Bell validity through two matrix inequalities};
\emph{An exact obstruction to the quadratic half-line completion},
and the accepted ramified order-49 data.
\href{https://github.com/anubhav70747/q222-finite-sdp/tree/16b37f6deb198fa48134145b4d12b7cc7d37be36/extracted-lifts/canonical-completion}{Public repository, canonical-completion, commit 16b37f6}.
\bibitem{assembly} \emph{Q222 Polar Assembly},
\href{https://github.com/anubhav70747/q222-finite-sdp/tree/16b37f6deb198fa48134145b4d12b7cc7d37be36/extracted-lifts/canonical-completion/full-polar}{same commit, full-polar}.
\bibitem{companion} Reconstructed R21 companion,
\href{https://github.com/anubhav70747/q222-finite-sdp/tree/16b37f6deb198fa48134145b4d12b7cc7d37be36/latex}{same commit, latex},
Chapters 3--6; in particular TF12, SG2, SG7, and CC1--CC6.
\end{thebibliography}

\clearpage
\appendix
\section{The complete integer moment certificate}\label{app:data}
All matrices below are used as integers in (5)--(10).
The machine-readable file also records the common normalization
$N/(15\cdot2^{60})$; multiplying a dual certificate by a positive
constant does not change its sign or its cancellations.

Positive definiteness can be checked without numerical eigenvalues.
For any matrix $R$ in (5), use the rational recurrence
\[
 p_j=R_{jj}-\sum_{k<j}L_{jk}^2p_k,\qquad
 L_{ij}=\frac{R_{ij}-\sum_{k<j}L_{ik}p_kL_{jk}}{p_j},
 \quad L_{jj}=1\quad(i>j).
\]
It gives $R=L\operatorname{diag}(p_j)L^T$. The exact floors of the
smallest pivots are
\[
\begin{array}{c|rrrr}
 R&\mathsf U_G&\mathsf V_G&\mathsf U_H&\mathsf V_H\\\hline
 \left\lfloor\min_jp_j\right\rfloor&794&12397&67408&1527738.
\end{array}
\]
Thus all four matrices are positive definite. The independent checker
instead uses fraction-free leading principal determinants and obtains
the same pivot bounds. The shared upper-right blocks in each pair
$N_{G,k},N_{H,k}$ make (6) directly visible.

\input{moments_g.tex}
\clearpage
\section*{The five $H$ moments}
\input{moments_h.tex}
\end{document}
